\begin{lexlean}{IdsSL}
\useglossary{lexlean.std.nat@1.2.0}
\useglossary{spec.titles@1.0.0}
\importmodule{Rows}
\title{Stdlib conformance claims}
\begin{semanticmodule}
definition{
  name: sl01,
  parameters: [],
  result: Rows::IdRow,
  body: (record Rows::IdRow {id := "SL-01", level := "build", suite := "stdlib", statement := "The standard library's `Text` module is language-1.2 source whose seven members over `Bytes` (`startsWith`, `endsWith`, `trim`, `lines`, `splitOn`, `joinWith`, `find`) are production roots eligible for `rust-std` with exactly their declared effects; each member's lowered program under the calculus evaluator and its rendered `rust-std` package under the pinned rustc agree with Rust's `starts_with`, `ends_with`, `trim_ascii`, `lines`, `split`, `join`, and `find` on a seeded corpus; a planted defect in a member is detected; and the certified packages the production export publishes for it are the committed ones.", sections := "§17.22"}),
  executable: true
}

definition{
  name: sl02,
  parameters: [],
  result: Rows::IdRow,
  body: (record Rows::IdRow {id := "SL-02", level := "build", suite := "stdlib", statement := "The standard library's `Sha256` module is language-1.2 source stating FIPS 180-4 SHA-256 over `UInt32` bit operations and `iterate` (modular addition by widening to `UInt64`, rotation as two shifts, padding read by position, the message schedule as a sixteen-word window, sixty-four rounds over the round-constant list), whose roots `sha256 : Bytes -> Digest` and `sha256Hex : Bytes -> Bytes` are production roots eligible for `rust-std` with exactly their declared effects; the lowered program under the calculus evaluator and the rendered `rust-std` package under the pinned rustc compute the FIPS 180-4 example digests and every NIST CAVP SHA-256 short- and long-message vector, and agree with the `sha2` crate on a seeded corpus whose lengths cover every padding boundary, a disagreement failing closed and never resolved by adopting the oracle's digest; a planted one-bit error in the round function is detected; and the certified packages the production export publishes for the two roots are the committed ones.", sections := "§17.22"}),
  executable: true
}

definition{
  name: sl05,
  parameters: [],
  result: Rows::IdRow,
  body: (record Rows::IdRow {id := "SL-05", level := "build", suite := "stdlib", statement := "The compiler computes its content identities with the standard library's certified `Sha256` package instead of the `sha2` crate: the shipped `lexlean` crate links the committed production export of `LexLeanStd.Sha256` through an in-crate link that resolves to that file and nothing else, its manifest has no `sha2` dependency, and `Sha256Digest::of`, `FramedHasher`, and `tree_digest` compute through the package; through the compiler's own API the digest of every example NIST states (the examples of FIPS 180-4 and the million-`a` message of FIPS 180-2) and every NIST CAVP SHA-256 short- and long-message vector, each cited by a `some-true` row of `model/authorities.toml` that vendors it or states its digest,, of every message length from 0 to 300 bytes and around each power of two up to 4 MiB, and of a seeded corpus agrees with the authority's digest and with the `sha2` crate as an oracle that can falsify and never establish, and the frames and tree rows agree with an independent recomputation of the §21.1 and §11.5 encodings, a disagreement failing closed and never resolved by adopting the oracle's digest; and a planted bit of the first round constant in a built copy of the package, a link that differs from the export by one byte, a `sha2` dependency added to the manifest, a message truncated by one byte, and a frame length written little-endian are each detected.", sections := "§17.22"}),
  executable: true
}
\end{semanticmodule}
\end{lexlean}
